% Supplementary Material of the RSE article: display details of the main figures. These details
% stood in the legends of manuscript/rse/fig_*.tex until the legends were shortened to fit the
% journal's word limit (October 2026); they are kept here with their substance unchanged, one
% paragraph per figure. The opening figure (fig_thennow) lost nothing and has no paragraph.
% Every number is a \metric{} macro; the article's labels are read through xr-hyper (prefix M-).

\subsection*{Display details of the main figures}
How the figures of the main article are drawn, and the values their legends leave out. Figure and section numbers refer to the main article.

\emph{Study area (figure~\ref{M-fig:rse_study_area} of the article).} Each sector is named on the map. The vermilion test blocks hold the ICESat-2 points used to test the elevation model. The emerged seabed (tan) is the seabed that was water in the Landsat extent of the first period (the \metric{params.highstand_label} high stand) and not in that of the second (the \metric{params.lowstand_label} low stand). The international boundary of panel a, and the land, lakes and boundaries of panel b, are from Natural Earth. Panel a uses the Lambert azimuthal equal-area projection centred at \metric{params.laea_lat0_deg}\textdegree{}N, \metric{params.laea_lon0_deg}\textdegree{}E, with a graticule every \metric{params.fig_graticule_deg}\textdegree{} and a scale bar of \metric{params.fig_scale_bar_km}~km.

\emph{The hindcast audit on one example (figure~\ref{M-fig:rse_method}).} Level bins are \metric{params.bin_width_m}~m wide. The acquisitions of the bin are split by their order into the two interleaved halves. The water frequency of panel a is the share of the valid acquisitions in which a cell was classified as water, and the observed shoreline is the edge of the cells whose frequency is at least \metric{params.contour_level} and that connect to the open sea. The offset $c$ of panel b was fitted to the flooded area on training blocks; a cell is flooded when a path from the open sea reaches it without crossing ground higher than the level. In panel c, the distances between predicted and observed shoreline cells pool the distances from both shorelines. In panel d, GEBCO\_2026 exceeds $Q$ at \metric{h1.n_exceeding.shape.gebco2026.interfluve} of the \metric{h1.n_evaluable.shape.gebco2026.interfluve} levels. Scale bars are \metric{params.rse_fig_method_scale_km}~km.

\emph{Emergence dates (figure~\ref{M-fig:rse_emergence}).} The \metric{params.dem_res_m}~m pixels of panel a are shown on a \metric{params.fig_map_res_m}~m display grid. Dates are coloured only on seabed that was water in \metric{params.highstand_label} and land in \metric{params.lowstand_label}, because elsewhere the elevation model is an extrapolation and a date would not mark an emergence. The graticule is drawn every \metric{params.fig_graticule_deg}\textdegree{}, and the scale bar is \metric{params.fig_scale_bar_km}~km.

\emph{Accuracy of the waterline elevation model (figure~\ref{M-fig:rse_dem_accuracy}).} The elevation model of panel a, at \metric{params.dem_res_m}~m, is shown on a \metric{params.fig_map_res_m}~m display grid. The ICESat-2 ATL08 height of panel b is the orthometric height plus the datum constant of \metric{h2.atl08.datum_constant_m}~m fitted on training blocks. The NMAD of panels b and d is \metric{params.nmad_scale} times the median absolute deviation from the median. The Kizlyar--Dagestan and Volga--Ural interfluve coasts hold \metric{h2.atl08.n_blocks.west} test block each, so no interval can be estimated for them.

\emph{Grid error against observation error (figure~\ref{M-fig:rse_dq}).} The degraded-elevation-model control drawn in grey is the waterline elevation model put into the grid's vertical frame, rounded to whole metres where the grid stores whole metres, averaged to the grid's cell size and audited exactly as the grid. Its ratio $D_\mathrm{control}/Q$ is shown at the levels where both the grid and the control can be measured. The shading marks displacements within the noise floor.

\emph{Shorelines of the five grids (figure~\ref{M-fig:rse_shorelines}).} The level bin was chosen by a fixed rule: in the three coasts of panels c--e, it is the late-fall bin at which the five grids' displacements are jointly closest to their medians over all levels. The observed shoreline cells are the \metric{params.stage2_res_m}~m waterline cells from which the displacement is measured. The grid heights of panel b are those the audit uses; the transect runs from the sea (left) to the land (right), and each grid's ticks, drawn in the grid's colour, mark where its predicted shoreline crosses the transect at the level. $Q$ in panels c--e is the \metric{params.err_pct_hi}th percentile of the distances between the shorelines of two interleaved halves of the acquisitions. The zooms have a graticule every \metric{params.erl_fig_graticule_zoom_deg}\textdegree{}, and zooms in one row share the scale of the \metric{params.erl_fig_scale_inset_km}~km bar.

\emph{Flooded area and shoreline position (figure~\ref{M-fig:rse_area}).} The area bias of panel a counts errors on cells of known state outside the observed area. The offsets fitted to flooded area on all periods' training blocks, used in panel a and for the open markers of panel b, are \metric{rev.w.s.oos.c_all_m_min} to \metric{rev.w.s.oos.c_all_m_max}~m; those fitted on the rising limb and early fall only (filled markers of panel b) are \metric{rev.w.s.oos.c_e12_m_min} to \metric{rev.w.s.oos.c_e12_m_max}~m. In panel c, the planar field lowers the summed absolute area error by \metric{rev.w.v.planar.reduction_pct_min}--\metric{rev.w.v.planar.reduction_pct_max}\% against one offset, while the median shoreline displacement rises from \metric{rev.w.v.planar.gebco2026.before_km} to \metric{rev.w.v.planar.gebco2026.after_km}~km for GEBCO\_2026 and from \metric{rev.w.v.planar.etopo1.before_km} to \metric{rev.w.v.planar.etopo1.after_km}~km for ETOPO1.

\emph{Great Salt Lake (figure~\ref{M-fig:rse_gsl}).} $Q$ is the \metric{params.err_pct_hi}th percentile of the distances between the shorelines of two interleaved halves of the acquisitions. The lake-bed control is built from USGS single-beam echo soundings below the sounded top and from digitised map contours and the National Elevation Dataset above it; its offset is \metric{rev2.g.lakebed.all.north.shape_offset_m}~m in the north arm and \metric{rev2.g.lakebed.all.south.shape_offset_m}~m in the south arm. The sounded levels (shaded columns of panels a and b) are bins \metric{rev2.g.lakebed.sounded.north.shape.bins} in the north arm (\metric{rev2.g.lakebed.coverage.north.sounded.level_low_navd88_m}--\metric{rev2.g.lakebed.coverage.north.sounded.level_high_navd88_m}~m) and \metric{rev2.g.lakebed.sounded.south.shape.bins} in the south arm (\metric{rev2.g.lakebed.coverage.south.sounded.level_low_navd88_m}--\metric{rev2.g.lakebed.coverage.south.sounded.level_high_navd88_m}~m). With the offset fitted over all its levels, the lake-bed control exceeds $Q$ at \metric{rev2.g.lakebed.all.north.class.sounded.n_exceeding} of \metric{rev2.g.lakebed.all.north.class.sounded.n_evaluable} sounded levels in the north arm and \metric{rev2.g.lakebed.all.south.class.sounded.n_exceeding} of \metric{rev2.g.lakebed.all.south.class.sounded.n_evaluable} in the south arm. With its offset refitted on the sounded levels it fails the test in both arms, so the control itself misplaces the sounded shorelines. The dotted lines of panel d mark the gauge levels of February 2000, \metric{rev.g.gsl.north.gauge_feb2000_navd88_m}~m in the north arm and \metric{rev.g.gsl.south.gauge_feb2000_navd88_m}~m in the south arm.

\emph{Near-term emergence (figure~\ref{M-fig:rse_nearterm}).} The count of panel a is drawn as the most frequent count per display block; the scale bar is \metric{params.fig_scale_bar_km}~km. The shading of panel b is computed at the reported falls and joined by straight lines; at $D = \metric{h3.scenario.1_m}$~m the five grids give \metricsf{rev2.n.reref.area.main.E.1.domain.min_km2}{2}--\metricsf{rev2.n.reref.area.main.E.1.domain.max_km2}{3}~\kmsq{}. The grey band of panel c is clipped to the axis.
